% Extracción quirúrgica del propietario V2 05b: líneas 100--359.
% El resto permanece en este archivo como procedencia, pero no se publica.
\iffalse
% Macroestructura solenoidal estrictamente interna del continuo discreto.
% Este módulo no introduce ninguna ecuación clásica ni un problema exterior.

\chapter{Macrostructure interne de bilan, mémoire et télescopage}
\label{chap:pdfv2-sol-interna}

\begin{statusbox}
L'unique entrée de ce chapitre est la structure discrète du continu.
La branche \(\mathrm{sol}\) se construit avant toute évaluation scalaire
extérieure et avant toute formulation différentielle classique. Sa chaîne
complète est
\[
 \mathcal C_{\rm cont}^{\rm disc}
 \supset\mathcal C_{\rm cont}^{\rm sol}
 \xrightarrow{\operatorname{Res}_{\rm sol}}
 \mathcal G_{\rm sol}
 \xrightarrow{\operatorname{pr}_{X,{\rm sol}}}
 X_{\chi,{\rm sol}}
 \xrightarrow{\mathcal A_{\rm sol}}
 \mathfrak M_{\rm sol}
 \xrightarrow{\mathcal P_{\rm sol}}
 \mathcal C_{\rm sol}^{\rm HMT}.
\]
Chaque flèche conserve comme coordonnées l'histoire et l'appareil qui la
produit. Aucune image de cette chaîne n'est un quotient qui oublie sa source.
\end{statusbox}

\section{Domaine propre avant toute évaluation ultérieure}
\label{sec:pdfv2-sol-dominio}

Soit
\[
 x=(x_0\xrightarrow{e_0}x_1\xrightarrow{e_1}\cdots)
 \in\mathcal C_{\rm cont}^{\rm disc}
\]
une histoire survivante. Chaque arête \(e_j\) conserve ses deux feuillets,
son orientation, son résidu, son quotient et sa retenue entière. Soient
\(N_j^+(e_j)\) et \(N_j^-(e_j)\) les multiplicités des deux feuillets et soit
\(\Delta_{\gamma_j}\in\mathbb Z\) le coefficient orienté du chemin, avec la
retenue de la composition déjà incorporée. L'accroissement signé interne est
\begin{equation}\label{eq:pdfv2-sol-b}
 b_j=b(e_j):=
 \Delta_{\gamma_j}\bigl(N_j^+(e_j)-N_j^-(e_j)\bigr)\in\mathbb Z.
\end{equation}
Il ne s'obtient pas en évaluant une trajectoire extérieure : c'est une coordonnée du
ledger de l'arête elle-même.

Ses parties positive et négative se construisent par
\begin{equation}\label{eq:pdfv2-sol-bpm}
 b_j^+:=\frac{|b_j|+b_j}{2},
 \qquad
 b_j^-:=\frac{|b_j|-b_j}{2}.
\end{equation}
Par conséquent,
\begin{equation}\label{eq:pdfv2-sol-bpm-identidades}
 b_j=b_j^+-b_j^-,
 \qquad
 |b_j|=b_j^++b_j^-,
 \qquad
 b_j^+b_j^-=0.
\end{equation}

Les deux mémoires primitives et leurs deux lecteurs sont
\begin{align}
 M_0^+=M_0^-&=0,
 &M_j^+&=\sum_{0\leq i<j}b_i^+,
 &M_j^-&=\sum_{0\leq i<j}b_i^-;
 \label{eq:pdfv2-sol-memorias}\\
 D_j^0&:=M_j^+-M_j^-,
 &H_j^0&:=M_j^++M_j^-.
 \label{eq:pdfv2-sol-dh}
\end{align}
En conséquence,
\begin{equation}\label{eq:pdfv2-sol-dh-diferencias}
 D_{j+1}^0-D_j^0=b_j,
 \qquad
 H_{j+1}^0-H_j^0=|b_j|.
\end{equation}
La mémoire signée \(D^0\) et la mémoire totale \(H^0\) restent des
objets distincts ; conserver seulement leur différence finale détruirait la
généalogie.

La sous-structure \(\mathcal C_{\rm cont}^{\rm sol}\) se compose des histoires
pour lesquelles, à tout préfixe et à tout raffinement nonadique :
\begin{enumerate}
 \item les accroissements \eqref{eq:pdfv2-sol-b} et les deux mémoires
       \eqref{eq:pdfv2-sol-memorias} existent ;
 \item les cylindres fins se regroupent en fibres de neuf extensions avec la
       même loi de probabilité projective ;
 \item les projections de profondeur et de résolution définies ci-dessous sont
       compatibles avec les extensions d'histoire ;
 \item chaque ligne de pertes et chaque colonne terminale est finie à horizon
       fini ;
 \item tout préfixe admet au moins une extension compatible à tout
       horizon ultérieur.
\end{enumerate}
Ce sont des conditions d'appartenance à la branche, non des prémisses tacites
attribuées à une histoire arbitraire.

\fi
\section{Raffinement nonadique commun : inclusion, espérance et mémoire de bilan}
\label{sec:pdfv2-sol-e9j9}

Soit \(Y_j(x)\) l'ensemble fini des cylindres compatibles avec le préfixe
\(x_{\leq j}\), muni de la mesure projective \(\mu_j\). La troncature
\[
 \tau_{j+1,j}:Y_{j+1}(x)\longrightarrow Y_j(x)
\]
a neuf extensions admissibles au-dessus de chaque cylindre. On considère les
espaces finis
\[
 \mathscr H_j(x):=\ell^2\bigl(Y_j(x),\mu_j\bigr).
\]
L'inclusion et l'espérance nonadiques sont
\begin{align}
 J_{9,j}:\mathscr H_j&\longrightarrow\mathscr H_{j+1},
 &(J_{9,j}f)(y)&=f(\tau_{j+1,j}y),
 \label{eq:pdfv2-sol-j9}\\
 E_{9,j}:\mathscr H_{j+1}&\longrightarrow\mathscr H_j,
 &(E_{9,j}g)(z)&=\frac1{9}
   \sum_{\tau_{j+1,j}y=z}g(y).
 \label{eq:pdfv2-sol-e9}
\end{align}
La compatibilité projective de \(\mu_{j+1}\) et \(\mu_j\) donne
\begin{equation}\label{eq:pdfv2-sol-ej-projector}
 E_{9,j}J_{9,j}=I_{\mathscr H_j},
 \qquad
 J_{9,j}^*=E_{9,j},
 \qquad
 P_{9,j}:=J_{9,j}E_{9,j}=P_{9,j}^*=P_{9,j}^2.
\end{equation}
Le projecteur microscopique complémentaire est
\begin{equation}\label{eq:pdfv2-sol-qmicro-preliminar}
 Q_{9,j}:=I-P_{9,j}.
\end{equation}

Si \(b_f\) est l'accroissement sur la fibre fine et
\begin{equation}\label{eq:pdfv2-sol-bcoarse}
 b_c:=E_{9,j}b_f,
\end{equation}
la convexité élémentaire de \(|\cdot|\) construit, et ne postule pas, la retenue de
compensation
\begin{equation}\label{eq:pdfv2-sol-kappa}
 \kappa_j
 :=\frac{E_{9,j}|b_f|-|b_c|}{2}\geq0.
\end{equation}
En effet, en utilisant \(t^\pm=(|t|\pm t)/2\), on obtient
\begin{align}
 E_{9,j}b_f^+
 &=\frac{E_{9,j}|b_f|+E_{9,j}b_f}{2}
   =b_c^++\kappa_j,
 \label{eq:pdfv2-sol-jensen-plus}\\
 E_{9,j}b_f^-
 &=\frac{E_{9,j}|b_f|-E_{9,j}b_f}{2}
   =b_c^-+\kappa_j.
 \label{eq:pdfv2-sol-jensen-minus}
\end{align}
Les deux feuillets perdent la même quantité lors de la compression ; c'est pourquoi la différence
signée est conservée tandis que la mémoire totale enregistre la compensation :
\begin{equation}\label{eq:pdfv2-sol-jensen-dh}
 E_{9,j}(b_f^+-b_f^-)=b_c,
 \qquad
 E_{9,j}(b_f^++b_f^-)=|b_c|+2\kappa_j.
\end{equation}

\section{Factorisation orthogonale des degrés de profondeur et des couches de résolution dans \(\mathscr H_j^{\mathrm{depth}}\widehat\otimes\mathscr H_j^{\mathrm{res}}\)}
\label{sec:pdfv2-sol-micro-shell}

Pour empêcher qu'une même perte soit comptée à la fois comme profondeur et
comme résolution, l'espace de chaque niveau est typé comme
\begin{equation}\label{eq:pdfv2-sol-factor-hilbert}
 \mathscr H_j
 =\mathscr H_j^{\rm depth}\widehat\otimes
  \mathscr H_j^{\rm res}.
\end{equation}
Pour \(j\geq1\), le premier facteur porte
\[
 P_j^{\rm depth}:=J_{9,j-1}E_{9,j-1}
 \quad\text{sur }\mathscr H_j^{\rm depth},
\]
et l'on fixe \(P_0^{\rm depth}=I\). Le second facteur porte la partition
finie des shells héritée de la frontière des cylindres. Si
\(\sigma_j:Y_j^{\rm res}\to S_j\) est cette partition, ses opérateurs sont
\begin{align}
 J_j^{\rm res}f&=f\circ\sigma_j,
 &E_j^{\rm res}g(s)&=
 \frac1{\#\sigma_j^{-1}(s)}
 \sum_{\sigma_j(y)=s}g(y),
 \label{eq:pdfv2-sol-shell-je}\\
 P_j^{\rm res}&:=J_j^{\rm res}E_j^{\rm res}.
 \label{eq:pdfv2-sol-shell-p}
\end{align}

On définit les trois projecteurs mutuellement orthogonaux
\begin{align}
 P_j^{\rm term}
 &:=P_j^{\rm depth}\otimes P_j^{\rm res},
 \label{eq:pdfv2-sol-pterm}\\
 Q_j^{\rm micro}
 &:=(I-P_j^{\rm depth})\otimes I,
 \label{eq:pdfv2-sol-qmicro}\\
 Q_j^{\rm shell}
 &:=P_j^{\rm depth}\otimes(I-P_j^{\rm res}).
 \label{eq:pdfv2-sol-qshell}
\end{align}
Alors
\begin{equation}\label{eq:pdfv2-sol-ortogonalidad}
 Q_j^{\rm micro}Q_j^{\rm shell}=0,
 \qquad
 I=P_j^{\rm term}+Q_j^{\rm micro}+Q_j^{\rm shell},
\end{equation}
et, pour tout \(\xi\in\mathscr H_j\),
\begin{equation}\label{eq:pdfv2-sol-pitagoras}
 \|\xi\|^2
 =\|P_j^{\rm term}\xi\|^2
  +\|Q_j^{\rm micro}\xi\|^2
  +\|Q_j^{\rm shell}\xi\|^2.
\end{equation}
Le terme
\((I-P_j^{\rm depth})\otimes(I-P_j^{\rm res})\) n'apparaît pas une seconde fois : il est déjà
contenu exactement une fois dans \(Q_j^{\rm micro}\).

\section{Extensions de cylindre et ligne complète de pertes}
\label{sec:pdfv2-sol-gamma-loss}

Une extension de cylindres du niveau \(j\) au niveau \(j+1\) induit l'application
\begin{equation}\label{eq:pdfv2-sol-gamma}
 \Gamma_j:\mathscr H_j\longrightarrow\mathscr H_{j+1},
 \qquad
 \Gamma_j f=f\circ\tau_{j+1,j},
\end{equation}
avec toutes les étiquettes de feuillet, d'orientation, de retenue, de mémoire et de chemin
transportées dans la base des cylindres. Pour \(n>j\), on écrit
\begin{equation}\label{eq:pdfv2-sol-gamma-composition}
 \Gamma_{j:n}:=\Gamma_{n-1}\cdots\Gamma_j,
 \qquad
 \Gamma_{j:j}:=I.
\end{equation}

Sur \(\mathscr H_j\), les accroissements non négatifs définissent des opérateurs de
multiplication
\begin{equation}\label{eq:pdfv2-sol-loss-multipliers}
 B_j^+:=M_{\sqrt{b_j^+}},
 \qquad
 B_j^-:=M_{\sqrt{b_j^-}},
 \qquad
 K_j:=M_{\sqrt{2\kappa_j}}.
\end{equation}
Chaque racine est prise point par point sur l'ensemble fini des cylindres ; si une
coordonnée est nulle, son opérateur est nul. L'espace des pertes est la somme
orthogonale externe
\[
 \mathscr E_j
 :=\mathscr H_j^{(+)}\oplus\mathscr H_j^{(-)}
   \oplus\mathscr H_j^{(\kappa)}
   \oplus\mathscr H_j^{({\rm micro})}
   \oplus\mathscr H_j^{({\rm shell})},
\]
et la ligne complète est
\begin{equation}\label{eq:pdfv2-sol-loss-row}
 L_j:\mathscr H_j\longrightarrow\mathscr E_j,
 \qquad
 L_j\xi
 =\bigl(
 B_j^+\xi,B_j^-\xi,K_j\xi,
 Q_j^{\rm micro}\xi,Q_j^{\rm shell}\xi
 \bigr).
\end{equation}
Par conséquent,
\begin{equation}\label{eq:pdfv2-sol-loss-norm}
 \|L_j\xi\|^2
 =\|B_j^+\xi\|^2+\|B_j^-\xi\|^2+\|K_j\xi\|^2
  +\|Q_j^{\rm micro}\xi\|^2
  +\|Q_j^{\rm shell}\xi\|^2.
\end{equation}
Les parties positive, négative, retenue de compensation, microstructure et shell
apparaissent chacune une fois. Elles ne fusionnent pas en une constante anonyme et ne se
répètent pas dans deux termes de la somme.

\section{Terminal d'histoire complète}
\label{sec:pdfv2-sol-terminal}

Soit \(\operatorname{Hist}_{J}^{\rm sol}\) l'ensemble des histoires
survivantes jusqu'à l'horizon \(J\), chacune avec son paquet de treize
coordonnées. L'espace terminal est défini comme l'espace libre de ces
registres complets :
\begin{equation}\label{eq:pdfv2-sol-terminal-space}
 \mathscr T_J^{\rm sol}
 :=\ell^2\!\left(
 \left\{(h,D_{{\rm sol},J}(h)):
 h\in\operatorname{Hist}_{J}^{\rm sol}\right\}
 \right).
\end{equation}
L'évaluation terminale
\begin{equation}\label{eq:pdfv2-sol-terminal-evaluation}
 E_J^{\rm term}:\mathscr H_J\longrightarrow\mathscr T_J^{\rm sol}
\end{equation}
envoie chaque vecteur de base de cylindre au registre complet de l'histoire qui
l'engendre et se prolonge linéairement. En particulier, le terminal retient
fibre, relations, nombres, orientation, retenue, mémoire, frontière, chemin,
multisection, survie et lecteur. Il n'est pas remplacé par la dernière valeur
visible du chemin.

\section{Colonne d'observation, Gram et télescopage}
\label{sec:pdfv2-sol-gram}

Pour \(0\leq j\leq J\), la colonne future complète est
\begin{equation}\label{eq:pdfv2-sol-g-column}
 \begin{aligned}
 G_{j;J}:\mathscr H_j&\longrightarrow
 \mathscr T_J^{\rm sol}\oplus
 \bigoplus_{n=j}^{J-1}\mathscr E_n,\\
 G_{j;J}\xi
 &:={}
 E_J^{\rm term}\Gamma_{j:J}\xi
 \ \oplus\ 
 \bigoplus_{n=j}^{J-1}L_n\Gamma_{j:n}\xi.
 \end{aligned}
\end{equation}
Pour \(j=J\), la somme des pertes est vide et
\(G_{J;J}=E_J^{\rm term}\). Le Gram terminal est
\begin{equation}\label{eq:pdfv2-sol-u-gram}
 U_{j;J}:=G_{j;J}^*G_{j;J}\succeq0.
\end{equation}

En réordonnant la somme orthogonale du codomaine, on obtient l'identité littérale
\begin{equation}\label{eq:pdfv2-sol-g-recursion}
 G_{j;J}\xi
 \simeq
 \bigl(L_j\xi,\,G_{j+1;J}\Gamma_j\xi\bigr).
\end{equation}
En conséquence,
\begin{equation}\label{eq:pdfv2-sol-stein}
 \boxed{U_{j;J}
 =L_j^*L_j+\Gamma_j^*U_{j+1;J}\Gamma_j}
\end{equation}
et
\begin{equation}\label{eq:pdfv2-sol-stein-difference}
 U_{j;J}-\Gamma_j^*U_{j+1;J}\Gamma_j
 =L_j^*L_j\succeq0.
\end{equation}
Il n'y a pas de mutation de signe : le terme de perte est du côté positif
de l'identité parce qu'il provient d'une somme de carrés.

En itérant \eqref{eq:pdfv2-sol-stein}, on obtient le télescopage exact
\begin{equation}\label{eq:pdfv2-sol-telescopia-operatorial}
 U_{j;J}
 =\Gamma_{j:J}^*(E_J^{\rm term})^*E_J^{\rm term}\Gamma_{j:J}
  +\sum_{n=j}^{J-1}
   \Gamma_{j:n}^*L_n^*L_n\Gamma_{j:n},
\end{equation}
ou, sous forme quadratique,
\begin{equation}\label{eq:pdfv2-sol-telescopia-normas}
 \langle U_{j;J}\xi,\xi\rangle
 =\|E_J^{\rm term}\Gamma_{j:J}\xi\|^2
  +\sum_{n=j}^{J-1}\|L_n\Gamma_{j:n}\xi\|^2.
\end{equation}
Le premier terme conserve le futur terminal complet ; la somme conserve,
niveau par niveau, toutes les pertes et toutes les retenues qui y ont conduit.

\iffalse
\section{Les treize coordonnées de la restriction}
\label{sec:pdfv2-sol-trece}

Pour un préfixe \(x_{\leq j}\), on définit
\begin{equation}\label{eq:pdfv2-sol-d-package}
 \begin{aligned}
 D_{{\rm sol},j}(x):=\bigl(&
 \mathcal F_{{\rm sol},j},
 \operatorname{Ext}_{\Gamma,{\rm sol},j},
 \operatorname{Rel}_{{\rm sol},j},
 \mathcal N_{{\rm sol},j},
 \operatorname{or}_{{\rm sol},j},
 \operatorname{Carr}_{{\rm sol},j},
 \operatorname{Mem}_{{\rm sol},j},\\
 &\partial_{{\rm sol},j},
 \gamma_{{\rm sol},j},
 \operatorname{Msect}_{{\rm sol},j},
 \operatorname{Surv}_{{\rm sol},j},
 \operatorname{Term}_{{\rm sol},j},
 \mathcal L_{{\rm sol},j}\bigr).
 \end{aligned}
\end{equation}
Ses composantes sont les suivantes.
\begin{enumerate}
 \item \(\mathcal F_{{\rm sol},j}\) contient les cylindres \(Y_j\), leur mesure
       projective, \(\mathscr H_j^{\rm depth}\),
       \(\mathscr H_j^{\rm res}\), leurs produits tensoriels et les cinq
       espaces de pertes.
 \item \(\operatorname{Ext}_{\Gamma,{\rm sol},j}\) contient
       \(J_{9,j},E_{9,j},\Gamma_j,\Gamma_{j:n}\), les inclusions de shell et
       toutes les applications de troncature et de raffinement.
 \item \(\operatorname{Rel}_{{\rm sol},j}\) contient
       \eqref{eq:pdfv2-sol-bpm-identidades},
       \eqref{eq:pdfv2-sol-ej-projector},
       \eqref{eq:pdfv2-sol-jensen-plus}--
       \eqref{eq:pdfv2-sol-jensen-minus},
       \eqref{eq:pdfv2-sol-ortogonalidad} et
       \eqref{eq:pdfv2-sol-stein}.
 \item \(\mathcal N_{{\rm sol},j}\) conserve
       \(j,J,9,b_j,b_j^+,b_j^-,\kappa_j\), les multiplicités de feuillet, les tailles
       des fibres et les dimensions des shells.
 \item \(\operatorname{or}_{{\rm sol},j}\) conserve
       \(\Delta_{\gamma_j}\), l'ordre du chemin et l'action d'inversion
       \(b_j\mapsto-b_j\), qui échange \(b_j^+\) et \(b_j^-\) sans changer
       \(|b_j|\) ni \(\kappa_j\).
 \item \(\operatorname{Carr}_{{\rm sol},j}\) conserve la retenue TPK entière,
       la retenue de composition et la retenue de compensation \(\kappa_j\) comme
       coordonnées différentes.
 \item \(\operatorname{Mem}_{{\rm sol},j}\) est
       \((M_j^+,M_j^-,D_j^0,H_j^0)\) avec le ledger ordonné de tous les
       accroissements antérieurs.
 \item \(\partial_{{\rm sol},j}\) contient la frontière orientée des cylindres,
       les accroissements \(b_j\), la ligne \(L_j\) et les différences de Stein.
 \item \(\gamma_{{\rm sol},j}\) est le chemin complet
       \(e_0e_1\cdots e_{j-1}\), et non seulement ses extrémités.
 \item \(\operatorname{Msect}_{{\rm sol},j}\) contient toutes les extensions
       nonadiques et de shell compatibles ; il ne choisit pas un prolongement.
 \item \(\operatorname{Surv}_{{\rm sol},j}\) enregistre les horizons
       \(J\geq j\) pour lesquels il existe \(G_{j;J}\), \(U_{j;J}\) et
       des terminaux compatibles.
 \item \(\operatorname{Term}_{{\rm sol},j}\) contient
       \(E_J^{\rm term}\Gamma_{j:J}\), le registre complet d'histoire et
       les treize coordonnées à chaque horizon.
 \item \(\mathcal L_{{\rm sol},j}\) publie en interne
       \((b^\pm,M^\pm,D^0,H^0,\kappa,Q^{\rm micro},Q^{\rm shell},
       L,G,U)\) et ses identités.
\end{enumerate}

Pour une extension \(u:x_j\to x_{j'}\), le transport coordonné
\(u_{{\rm sol},j',j}\) satisfait
\begin{equation}\label{eq:pdfv2-sol-naturality-d}
 u_{{\rm sol},j',j}D_{{\rm sol},j'}(x_{j'})
 =D_{{\rm sol},j}(\tau_{j',j}x_{j'}).
\end{equation}
L'égalité inclut les cinq composantes de \(L\), le terminal complet et
la composition de \(\Gamma\) ; ce n'est pas une égalité des seuls cardinaux.

\section{Graphe de restriction et objet dépendant}
\label{sec:pdfv2-sol-graphs}

On définit
\begin{equation}\label{eq:pdfv2-sol-g-graph}
 \mathcal G_{\rm sol}:=\operatorname{Graph}(D_{\rm sol}),
 \qquad
 \operatorname{Res}_{\rm sol}(x):=(x,D_{\rm sol}x).
\end{equation}
La première projection est inverse sur le graphe :
\begin{equation}\label{eq:pdfv2-sol-res-inverse}
 \pi_1\operatorname{Res}_{\rm sol}=I,
 \qquad
 \operatorname{Res}_{\rm sol}\pi_1=I_{\mathcal G_{\rm sol}}.
\end{equation}

Soit \(\chi_{\rm sol}(x)\) la multisection complète formée de tous les
cylindres, raffinements, shells, lignes de pertes et terminaux compatibles
avec \(x\). Alors
\begin{equation}\label{eq:pdfv2-sol-xchi}
 X_{\chi,{\rm sol}}
 :=\{(\chi_{\rm sol}(x),x,D_{\rm sol}x):
 x\in\mathcal C_{\rm cont}^{\rm sol}\},
\end{equation}
et
\begin{equation}\label{eq:pdfv2-sol-prx}
 \operatorname{pr}_{X,{\rm sol}}(x,D_{\rm sol}x)
 :=(\chi_{\rm sol}(x),x,D_{\rm sol}x).
\end{equation}
De nouveau, la projection qui conserve \((x,D_{\rm sol}x)\) est inverse sur
cette image dépendante.

\section{Assembleur et publicateur internes}
\label{sec:pdfv2-sol-assembler-publisher}

L'espace ambiant \(\mathscr M_{\rm sol}^{\rm amb}\) se compose de systèmes
compatibles
\[
 \bigl(
 \{b,b^\pm,M^\pm,D^0,H^0,\kappa\},
 \{E_9,J_9,P^{\rm term},Q^{\rm micro},Q^{\rm shell}\},
 \{\Gamma,L,E^{\rm term},G,U\}
 \bigr)
\]
avec toutes les identités précédentes. L'assembleur brut est
\begin{equation}\label{eq:pdfv2-sol-assembler}
 \begin{aligned}
 a_{\rm sol}:X_{\chi,{\rm sol}}&\longrightarrow
 \mathscr M_{\rm sol}^{\rm amb},\\
 a_{\rm sol}(\chi(x),x,Dx)&:=
 \bigl(
 \{b,b^\pm,M^\pm,D^0,H^0,\kappa\}_x,
 \{E_9,J_9,P,Q\}_x,
 \{\Gamma,L,E^{\rm term},G,U\}_x
 \bigr).
 \end{aligned}
\end{equation}
La macrostructure conserve son antécédent par
\begin{equation}\label{eq:pdfv2-sol-macro-graph}
 \mathfrak M_{\rm sol}:=\operatorname{Graph}(a_{\rm sol}),
 \qquad
 \mathcal A_{\rm sol}(z):=(z,a_{\rm sol}z).
\end{equation}

Soit \(\mathscr Y_{\rm sol}^{\rm int}\) le type des certificats formés par
les identités
\begin{equation}\label{eq:pdfv2-sol-certificate-list}
 \begin{gathered}
 b=b^+-b^-,\qquad |b|=b^++b^-,\\
 E_9J_9=I,\qquad J_9E_9=P_9,\\
 E_9b_f^\pm=b_c^\pm+\kappa,\qquad\kappa\geq0,\\
 I=P^{\rm term}+Q^{\rm micro}+Q^{\rm shell},
 \qquad Q^{\rm micro}Q^{\rm shell}=0,\\
 U_{j;J}=L_j^*L_j+\Gamma_j^*U_{j+1;J}\Gamma_j,\\
 \langle U_{j;J}\xi,\xi\rangle
 =\|E_J^{\rm term}\Gamma_{j:J}\xi\|^2
  +\sum_{n=j}^{J-1}\|L_n\Gamma_{j:n}\xi\|^2,
 \end{gathered}
\end{equation}
avec la naturalité de toutes les flèches. Le publicateur brut
\begin{equation}\label{eq:pdfv2-sol-publisher}
 p_{\rm sol}^{\rm raw}:\mathfrak M_{\rm sol}
 \longrightarrow\mathscr Y_{\rm sol}^{\rm int}
\end{equation}
publie ce certificat sans supprimer le macro-objet. Enfin,
\begin{equation}\label{eq:pdfv2-sol-output-graph}
 \mathcal C_{\rm sol}^{\rm HMT}
 :=\operatorname{Graph}(p_{\rm sol}^{\rm raw}),
 \qquad
 \mathcal P_{\rm sol}(m):=(m,p_{\rm sol}^{\rm raw}m).
\end{equation}

\section{Théorème interne de génération solénoïdale}
\label{sec:pdfv2-sol-theorem}

\begin{theorem}[Génération interne de la macrostructure de bilan]
\label{thm:pdfv2-sol-internal-generation}
Pour tout \(x\in\mathcal C_{\rm cont}^{\rm sol}\), la chaîne
\[
 \mathcal C_{\rm cont}^{\rm sol}
 \xrightarrow{\operatorname{Res}_{\rm sol}}
 \mathcal G_{\rm sol}
 \xrightarrow{\operatorname{pr}_{X,{\rm sol}}}
 X_{\chi,{\rm sol}}
 \xrightarrow{\mathcal A_{\rm sol}}
 \mathfrak M_{\rm sol}
 \xrightarrow{\mathcal P_{\rm sol}}
 \mathcal C_{\rm sol}^{\rm HMT}
\]
construit naturellement une macrostructure qui contient les deux mémoires,
la retenue de compensation, la séparation orthogonale micro/shell, la ligne complète
de pertes, la colonne future, le terminal d'histoire complète, le Gram et le
télescopage. Les treize coordonnées du continu peuvent être récupérées par
projections successives depuis n'importe quel point de la chaîne.
\end{theorem}

\begin{proof}
Les identités \eqref{eq:pdfv2-sol-bpm-identidades} et
\eqref{eq:pdfv2-sol-dh-diferencias} s'obtiennent directement à partir de la
décomposition positive et négative de l'accroissement entier ; les
mémoires ne sont donc pas des données ajoutées. Les formules
\eqref{eq:pdfv2-sol-j9}--\eqref{eq:pdfv2-sol-e9}, avec la mesure
projective, prouvent \(E_9J_9=I\), \(J_9^*=E_9\) et que \(P_9\) est un
projecteur orthogonal.

L'inégalité triangulaire donne
\(E_9|b_f|\geq|E_9b_f|=|b_c|\), de sorte que
\(\kappa\geq0\). Substituer
\(b^\pm=(|b|\pm b)/2\) prouve
\eqref{eq:pdfv2-sol-jensen-plus} et
\eqref{eq:pdfv2-sol-jensen-minus} sans hypothèse supplémentaire.

Les définitions tensorielles
\eqref{eq:pdfv2-sol-pterm}--\eqref{eq:pdfv2-sol-qshell} prouvent que
\(P^{\rm term}\), \(Q^{\rm micro}\) et \(Q^{\rm shell}\) sont orthogonaux et
ont pour somme l'identité. Par conséquent, la ligne \(L_j\) compte chaque secteur une
seule fois.

La définition \eqref{eq:pdfv2-sol-g-column} sépare orthogonalement le terminal
et les pertes de chaque niveau. Séparer le premier terme de pertes produit
\eqref{eq:pdfv2-sol-g-recursion}. Prendre l'adjoint et composer prouve l'identité
de Stein \eqref{eq:pdfv2-sol-stein} ; l'itérer prouve
\eqref{eq:pdfv2-sol-telescopia-operatorial} et
\eqref{eq:pdfv2-sol-telescopia-normas}. Tout terme est une norme au carré,
de sorte que le signe opposé ne peut apparaître.

La compatibilité de la troncature des cylindres, de l'espérance, des projections et des
extensions \(\Gamma\) donne la naturalité
\eqref{eq:pdfv2-sol-naturality-d}. Enfin,
\(\mathcal G_{\rm sol}\), \(\mathfrak M_{\rm sol}\) et
\(\mathcal C_{\rm sol}^{\rm HMT}\) sont des graphes, et
\(X_{\chi,{\rm sol}}\) est un objet dépendant qui conserve explicitement
son antécédent. La première projection de chaque étape récupère l'étape
précédente. Ainsi aucune composition n'efface l'histoire, la mémoire, la multisection,
le terminal ni le lecteur.
\end{proof}

\section{Provenance probatoire et falsificateurs}
\label{sec:pdfv2-sol-provenance-falsifiers}

\paragraph{Résultat récupéré.}
Appartiennent à l'architecture déjà construite l'accroissement signé avec retenue,
la séparation \(b=b^+-b^-\), les mémoires \(M^\pm,D^0,H^0\), le couple
\(E_9,J_9\), la retenue de Jensen \(\kappa\), l'extension des cylindres, la
ligne de pertes, le terminal, le Gram, Stein et le télescopage.

\paragraph{Formalisation nouvelle.}
Sont des formalisations explicites de cette architecture le typage préalable à toute
évaluation extérieure, la factorisation
\(\mathscr H^{\rm depth}\widehat\otimes\mathscr H^{\rm res}\), la séparation
orthogonale qui empêche le double comptage, le terminal libre sur des histoires à
treize coordonnées et les quatre mises en paquet non ablatives par graphes.

\begin{criterion}[Falsificateurs internes de la branche \(\mathrm{sol}\)]
\label{crit:pdfv2-sol-falsifiers}
La construction échoue si l'un quelconque des faits suivants se produit :
\begin{enumerate}
 \item \(b\neq b^+-b^-\) o \(|b|\neq b^++b^-\);
 \item \(E_9J_9\neq I\), la mesure n'est pas projective ou \(P_9\) n'est pas un
       projecteur ;
 \item \(\kappa<0\) ou
       \(E_9b_f^\pm=b_c^\pm+\kappa\) échouent ;
 \item \(Q^{\rm micro}Q^{\rm shell}\neq0\), ou un même secteur apparaît dans
       les deux termes de la ligne \(L\) ;
 \item l'une de \(b^+,b^-,\kappa,Q^{\rm micro},Q^{\rm shell}\) est omise ou
       comptée deux fois ;
 \item le terminal conserve seulement la valeur visible et non l'histoire avec ses
       treize coordonnées ;
 \item l'identité de Stein échoue ou le télescopage contient un signe
       contraire à une somme de carrés ;
 \item une extension ne commute pas avec \(E_9,J_9,L,G,U\) ;
 \item la multisection des continuations est remplacée par un prolongement
       choisi rétrospectivement ;
 \item l'une quelconque des treize coordonnées est omise.
\end{enumerate}
Dans le dernier cas, le diagnostic contraignant est : \emph{objet substitué ;
conclusion non transférable}.
\end{criterion}

\begin{warningbox}
Ce chapitre conclut uniquement la construction interne de
\(\mathcal C_{\rm sol}^{\rm HMT}\). Il ne contient pas d'équation différentielle
classique, de données initiales classiques, d'application à partir de données extérieures ni de
conclusion sur un problème extérieur. Ces entrelaceurs appartiennent à un autre
document et ne gouvernent pas l'existence de cette macrostructure.
\end{warningbox}
\fi
